\newpage 
\chapter{Numerical Integration}

The main idea of the numerical integration is to find a function \(g(x)\) which is easier to integrate, instead 
of integrating \(f(x)\) directly.

\section{Newton Cotes}

Given \(- \infty < a < b < \infty \) and \(f \in C[a,b]\) we want to approximate 

\[
    \int_a^b f(x) dx
\]

with use of equidistant points \(x_k = a(b - a)\frac{k}{n} \in [a, b]\) \(k\) being the number of points.
Hence, we use polynomial interpolation which gives us: 

\[
    \int_a^b g(x) dx = (b - a) \sum_{k = 0}^{n} \beta_k f(x_k)
\]

where \(\beta_k\) are the \emph{weights} given by our Newton interpolation algorithm.

It is also a preferred to divide the target interval into sub-intervals and then apply the approximation on those intervals to avoid 
using high degree polynomials.

\[
    \int_a^b f(x) dx = \sum_{i = 0}^{m - 1} \int_{a_i}^{b_i} g(x)
\]

with \(a = a_0 < b_0 = a_1 < b_1 ... a_{m - 1} < b_{m - 1} = b\).

\textbf{Example:}

\section{Lagrange Integration For \texorpdfstring{\([0, 1]\)}{}}

For the interval \([0,1]\) we get the equidistant points \(x_k = \frac{k}{n}\) giving 
us the interpolation 

\[
    g(x) = \sum_{k = 0}^{n} f(x_k) L_k (x),
\]

where 

\[
    L_k(x) = \prod_{i = 0, i \ne k} \frac{x - x_i}{x_k - x_i}
\]

is our \emph{lagrange polynomial}.

Using them gives us the approximation 

\[
    \int_1^0 g(x) = \sum_{k = 0}^{n}f(x_k) \int_0^1 L_k (x) dx = \sum_{k = 0}^n \beta_k f(x_k)
\]

with \(\beta_k = \int_0^1 L_k (x) dx\).

\textbf{Example:}

\subsection{Generalization For \texorpdfstring{\([a,b]\)}{}}

We can extend the previous idea to any real, closed interval 

\[
    \int_a^b f(x) dx = (b - a) \int_0^1 f(a + (b - a)y) dy \approx (b - a) \sum_{k = 0}^n \beta_k f \left(a + (b - a)\frac{k}{n}\right)
\]

\section{Summation For Error Mitigation}

\[
    \int_a^b f(x) dx = (b - a) \sum_{i = 0}^{m - 1} \int_{a_i}^{b_i} f(x) dx \approx \sum_{i = 0}^{m - 1}(b_i - a_i) \sum_{k = 0}^n \beta_k f(x_{k,i})
\]

with \(a = a_0 < b_0 = a_1 < b_1 ... a_{m - 1} < b_{m - 1} = b\). \(x_{k, i} = a_i + (b_i - a_i)\frac{k}{n}\).

For the case when the intervals have the same size then our intervals boundaries are given by 

\[
    h = \frac{b - a}{m} = b_i - a_i, \quad a_i = a + ih, \, b_i = a + (i + 1)h
\]

\textbf{Example:}

\section{Error Analysis}

\subsection{Newton Cotes}


Given \(f \in C^{2 + n} [0, h], \, I_n (f)\), then there is \(\xi \in [0, h]\) such that 

\[
    E_n(h) = \int_0^h f(x) dx - NC_n(h) = 
    \begin{cases}
        \frac{h}{n}^{n + 3} \int_0^n \frac{t \prod_{k = 0}^{n}(t - k)}{(n + 2)!} dt f^{(n + 2)}(\xi), & n \text{ even }\\
        \frac{h}{n}^{n + 2} \int_0^n \frac{\prod_{k = 0}^{n}(t - k)}{(n + 1)!} dt f^{(n + 1)}(\xi), & n \text{ odd }
    \end{cases}
\]

\subsection{Degree Of Exactness}

If \(I(f)\) is an integral approximation, which for all polynomials of degree \(\le l\) returns the exact solution, then 
we define \(l\) as the \emph{degree of exactness}.



\section{Gaussian Integration}





